\Legend
设放完后第 $i$ 列的总高度为 $b_i\ge a_i$。竖直支撑要求每列箱子连续。若第 $i$ 列放了箱子，则最高那个箱子的左右两侧必须有支撑，故 $b_i\le b_{i-1}$ 且 $b_i\le b_{i+1}$。因此，一段连续放过箱子的列必然等高，且它两端未放箱子的原货架高度都至少为这个高度。

于是一个方案可唯一分解为若干个互不相交的“填谷区间”。对于两个未放箱子的边界列 $u<v$，若中间各列均低于高度 $H$、而两边的原高度均至少为 $H$，就可以把 $u+1,\ldots,v-1$ 全部填到 $H$。记中间长度为 $L=v-u-1$，原高度和为 $S$，原高度最大值为 $M$，则可选的 $H$ 满足

$$
M < H \le \min(a_u,a_v), \qquad \text{箱数}=LH-S.
$$

把 $a_0,a_{n+1}$ 视为足够大的数。两边都高于区间内部的边界对只有 $O(n)$ 个：按右端点扫描时，单调栈每次弹出的元素和弹完后的栈顶各产生至多一对。实现中也可以用 $O(n^2)$ 扫描全部候选边界，遇到内部最大值不低于左端高度时停止；扫描量至多 $10^6$。

令 $dp[v][j]$ 表示处理到第 $v$ 个原货架且总共放了 $j$ 个箱子的方案数。可以跳过当前货架（权重为 $0$），也可以从边界 $u$ 跳到 $v$ 并填充两者之间的谷。固定一对边界后，允许的箱数构成公差为 $L$ 的等差数列。对每个余数类用滑动窗口，把该转移在 $O(m)$ 时间内完成，避免枚举高度后再对每个箱数循环。

答案为 $dp[n+1][m]$。时间复杂度 $O(n^2+nm)$，空间复杂度 $O(nm)$。
